\section{Denoising methods}

In this section, we present several denoising methods for a certain signal $x$ corrupted by noise to produce a noisy observation $y$.
Let $(\Omega, \mathcal{F}, \mathbb{P})$ be a probability space.
Let $X : \Omega \to \mathbb{R}^n$ and $Y : \Omega \to \mathbb{R}^n$ be random variables representing respectively the clean signal and the noisy observation.
We consider a corruption process $\mathcal{C}$ that maps the clean signal $X$ to the noisy observation $Y$ according to a certain noise model. This corruption process can be expressed as follows :

\begin{equation}
    Y | X=x \sim \mathcal{C}(x)
    \label{eq:noise-model}
\end{equation}

We will consider the residual additive noise $E$ defined as follows :

\begin{equation}
    E = Y - X
\end{equation}

\subsection{Self-supervised denoising}

\subsubsection{Denoising Auto-Encoders}

Whereas traditional auto-encoders are trained to reconstruct their input through an encoder-decoder structure, Denoising auto-encoders (DAEs) are trained to reconstruct clean data from corrupted versions of it.
They have been first introduced by \cite{vincent_extracting_2008} and have since been widely used in various applications such as image denoising, inpainting or other inverse problems.
The DAE objective is thus :

\begin{equation}
    \mathcal{L}_{DAE}(\theta) = \underset{\substack{X \sim p_X \\ Y | X=x \sim \mathcal{C}(x)}}{\mathbb{E}} \left[ \| X - f_{\theta}(Y) \|_2^2 \right]
\end{equation}
with $\mathcal{C}$ being a random corruption process.

\subsubsection{Score-based denoising}


\begin{definition}{Score function}
    % The score function of a random variable $X$ with $\mathcal{C}^1$ probability density function $p(x)$ is defined as the gradient of the log-density with respect to $x$ :
    \begin{equation}
        s_X(x) = \nabla_x \ln p(x)
    \end{equation}
\end{definition}

This function denotes the direction of maximum increase of the probability density function at point $x$.

\paragraph{Score matching estimation}

As first introduced by \cite{hyvarinen_estimation_2005}, score matching is a method to estimate parametric models of probability density functions without requiring the computation of the normalizing constant.
Given a parametric model score function $s(\cdot; \theta)$ with parameter $\theta$, and the data score function $s_X(\cdot)$ of observed data sampled from $X$, we want to minimize the distance between the data score function $s_X(\cdot)$ and the model score function $s_X(\cdot; \theta)$ using the following objective :

\begin{equation}
    \mathcal{L}_{SM}(\theta) = \mathbb{E}_X \left[ \| s_X(x) - s(x; \theta) \|_2^2 \right]
\end{equation}

Minimization of this objective does not require knowledge of the normalizing constant of $s(\cdot; \theta)$ as it only involves derivatives of the log-density.
They also show that we don't need to know the data score function $s_X(x)$ to optimize this objective as it can be rewritten as :

\begin{equation}
    \mathcal{L}_{SM}(\theta) = \mathbb{E}_X \left[ \mathop{Tr} \left( \nabla_x s(x; \theta) \right) + \frac{1}{2} \| s_X(x; \theta) \|_2^2 \right] + C
\end{equation}

This is very convenient as it avoids the need to estimate the data score function over the observed data, which is generally intractable.

\paragraph{Denoising Auto-Encoders and score function estimation}




DAEs and score matching principles have been linked together by \cite{vincent_connection_2011}. They built an objective similar to score matching but defined on clean/noisy image pairs with the corruption process $\mathcal{C}$ distribution being known.

\begin{equation}
    \mathcal{L}_{DSM}(\theta) = \underset{X \sim p_X \\ Y \sim \mathcal{C}(X)}{\mathbb{E}} \left[ \| s_X(Y, \theta) - s_{Y|X}(Y|X) \|_2^2 \right]
\end{equation}

They prove the specific case of Tweedie's formula for additive Gaussian noise (see Theorem \ref{theorem:tweedie_Gaussian}) which will be considered in the following.

\cite{alain_what_2014} further prove that when the Gaussian noise level $\sigma$ tends to zero, the DAE learns the score of the data distribution as follows :

\begin{equation}
    f_{\theta^*}(y) = y + \sigma^2 s_X(y) + \underset{\sigma \to 0}{o}(\sigma^2)
\end{equation}

This can be generalized to a wider range of exponential family distributions thanks to Tweedie's formula \cite{efron_tweedies_2011}. We retain the two following theorems that stand for additive Gaussian noise and Poisson noise respectively.

\begin{theorem}{Tweedie's formula for additive Gaussian noise}
    \label{theorem:tweedie_Gaussian}
    Let $X$ be a random variable and $Y = X + N$ be a noisy observation of $X$ corrupted by an additive Gaussian noise $N$ independent from $X$ with variance $\sigma^2$. The MMSE estimator of $X$ given $Y = y$ is given by :
    \begin{equation}
        \mathbb{E}\left[X | Y = y\right] = y + \sigma^2 \nabla_y \ln p(y)
    \end{equation}
\end{theorem}

\begin{proof}
    Bayes' rule gives us:
    \begin{equation*}
        p(x|y) = \frac{p(y|x)p(x)}{p(y)}
    \end{equation*}
    The likelihood $p(y|x)$ is given by the noise distribution:
    \begin{equation*}
        p(y|x) = \frac{1}{\sqrt{2 \pi \sigma^2}} \exp\left(-\frac{||y - x||^2}{2 \sigma^2}\right)
    \end{equation*}
    Thus we have:
    \begin{equation*}
        p(x|y) = \frac{1}{Z(y)} \exp\left(-\frac{||y - x||^2}{2 \sigma^2}\right) p(x)
    \end{equation*}
    with $Z(y)$ acting as a normalizing constant:
    \begin{equation*}
        Z(y) = \int p(x) \exp\left(-\frac{||y - x||^2}{2 \sigma^2}\right)dx
    \end{equation*}
    Then we compute the derivative of $p(y)$:
    \begin{align*}
        \nabla_y p(y) &= \frac{d}{dy} \int p(y|x)p(x) dx \\
        &= \int \frac{\partial}{\partial y} p(x) \exp\left(-\frac{||y - x||^2}{2 \sigma^2}\right) dx \\
        &= - \frac{1}{\sigma^2} \int (y - x) p(x) \phi_{\sigma}(y - x) dx \\
    \end{align*}
    where we denote $\phi_{\sigma}$ the Gaussian kernel of standard deviation $\sigma$.

    \begin{align*}
        \nabla_y \ln p(y) = \frac{p'(y)}{p(y)} &= - \frac{1}{\sigma^2} \left( \int (y - x) \frac{p(x) \phi_{\sigma}(y - x)}{p(y)} dx \right) \\
        &= - \frac{1}{\sigma^2} \left( y - \int x p(x|y) dx \right) \\
        &= - \frac{1}{\sigma^2} (y - \mathbb{E}[x|y])
    \end{align*}
\end{proof}

\begin{theorem}{Tweedie's formula for Poisson noise}
    Let $X$ be a random variable and $Y$ be a noisy observation of $X$ corrupted by Poisson noise. The MMSE estimator of $X$ given $Y = y$ is given by :
    \begin{equation}
        \mathbb{E}\left[X | Y = y\right] = (y + 1)\exp \left( \nabla_y \ln p(y)\right)
    \end{equation}
\end{theorem}

\begin{proof}
    The assumption of Poisson noise means that the likelihood $p(y|x)$ is given by:
    \begin{equation*}
        p(y|x) = \frac{(x)^{y} e^{-x}}{y!}
    \end{equation*}
    Then we compute the gradient of $p(y)$:
    \begin{align*}
        \exp(\ln \nabla_y p(y)) &= \frac{p(y+1)}{p(y)} \\
        &= \frac{1}{p(y)} \int p(x) \frac{x}{y+1}\frac{x^y}{y!}e^{-x} dx \\
        &= \frac{1}{y+1} \int x \frac{p(x) p(y|x)}{p(y)} dx \\
        &= \frac{1}{y+1} \mathbb{E}[x|y]
    \end{align*}
    Thus we have:
    \begin{equation*}
        \mathbb{E}[x|y] = (y+1) \exp(\ln \nabla_y p(y))
    \end{equation*}
\end{proof}


Given the previous proposition and theorems, one can build a link between denoising auto-encoders and score function estimation as follows :

\begin{equation}
    s_Y(y) = \frac{D_{\sigma}(y) - y}{\sigma^2} \quad \text{(Gaussian noise)}
\end{equation}
\begin{equation}
    s_Y(y) = \ln \left( \frac{D_{Poisson}(y)}{y + 1} \right) \quad \text{(Poisson noise)}
\end{equation}

We can highlight an intriguing application of \cite{kim_noise2score_nodate} who leverage this link to estimate the score function using a DAE trained under an additive Gaussian noise assumption, even when the actual noise model is not Gaussian. They then use the estimated score function to denoise images corrupted by other types of noise using the appropriate variant of Tweedie's formula.

\subsubsection{Stein's Unbiased Risk Estimator}